% GENERATED FILE. Edit canonical lessons, then run npm run generate:books.
% canonical-source-hash: fnv1a64:ebd48fa9101f35ce
% canonical-lessons: SA-C485-susrusa, SA-C485-upacarati, SA-C485-apyayayati, SA-C485-upasamah, SA-C485-rasayanam

\chapter{Nursing}
\label{ch:sa-nursing}
\hlchaptermodality{485}

\begin{chapteropening}
\textbf{By the end of this chapter:} \emph{I can say how I care for someone who is ill.}

Complete the last lesson of chapter 485: I can say how I care for someone who is ill.
\end{chapteropening}

\section[\texorpdfstring{\sk{शुश्रूषा} (śuśrūṣā)}{शुश्रूषा (śuśrūṣā)}]{\sk{शुश्रूषा} (śuśrūṣā) — nursing; dutiful service}

\label{lesson:SA-C485-susrusa}



\begin{quote}
Before the new one: say the Sanskrit for phlegm, then the Sanskrit for a wholesome diet.
\end{quote}

\subsection*{You'll want to know: \sk{शुश्रूषा}}
\textbf{\sk{शुश्रूषा}} — \emph{śuśrūṣā} — \textquotedblleft{}nursing; dutiful service\textquotedblright{}.

\begin{quote}
\sk{परिचारिकायाः} \sk{शुश्रूषया} \sk{आतुरः} \sk{स्वस्थः} \sk{भवति।}
\end{quote}

\begin{quote}
\emph{paricārikāyāḥ śuśrūṣayā āturaḥ svasthaḥ bhavati.}
\end{quote}

\begin{quote}
With the nurse's care the patient gets well.
\end{quote}

A feminine noun.

\begin{cousinweb}
\textbf{\sk{शुश्रूषा}} (\emph{śuśrūṣā}) is built on the desiderative of the root √śru, \textquotedblleft{}to hear\textquotedblright{}, the root of \textbf{\sk{शृणोति}}: MW's first sense is \textquotedblleft{}desire to hear\textquotedblright{}, then \textquotedblleft{}reverence, obedience, service\textquotedblright{}; \textquotedblleft{}nursing the sick\textquotedblright{} is the modern sense (MW).
\end{cousinweb}

\begin{grammarlens}[title={What you've built}]
One more word for describing what you lived and felt.
\end{grammarlens}

\subsection*{Your turn}
\begin{itemize}
\raggedright
  \item \emph{Say it:} \emph{śuśrūṣā}
  \item \emph{Say it:} \emph{śuśrūṣā}, once more
  \item \emph{Say it:} \emph{paricārikāyāḥ śuśrūṣayā āturaḥ svasthaḥ bhavati.} — With the nurse's care the patient gets well.
  \item \emph{Recall:} say the Sanskrit for phlegm, then the Sanskrit for a wholesome diet, then say \emph{śuśrūṣā} again
\end{itemize}

\begin{tcolorbox}[breakable,colback=teal!4,colframe=teal!35!black,title={Before you move on}]
What does \sk{शुश्रूषा} mean? (\textquotedblleft{}Nursing; dutiful service\textquotedblright{}.) Say it once more.
\end{tcolorbox}
\section[\texorpdfstring{\sk{उपचरति} (upacarati)}{उपचरति (upacarati)}]{\sk{उपचरति} (upacarati) — attends on}

\label{lesson:SA-C485-upacarati}



\begin{quote}
Before the new one: say the Sanskrit for a wholesome diet, then the Sanskrit for nursing.
\end{quote}

\subsection*{You'll want to know: \sk{उपचरति}}
\textbf{\sk{उपचरति}} — \emph{upacarati} — \textquotedblleft{}attends on\textquotedblright{}.

\begin{quote}
\sk{पुत्री} \sk{रुग्णां} \sk{मातरम्} \sk{उपचरति।}
\end{quote}

\begin{quote}
\emph{putrī rugṇāṃ mātaram upacarati.}
\end{quote}

\begin{quote}
The daughter looks after her sick mother.
\end{quote}

Past (she …): \sk{सा} \sk{मातरम्} \sk{उपचरितवती}, \emph{sā mātaram upacaritavatī}.

\begin{cousinweb}
\textbf{\sk{उपचरति}} (\emph{upacarati}) is built from \emph{upa}, \textquotedblleft{}near\textquotedblright{}, and the root √car, \textquotedblleft{}to move\textquotedblright{}; MW gives \textquotedblleft{}to wait upon, serve, attend\textquotedblright{} (MW).
\end{cousinweb}

\begin{grammarlens}[title={What you've built}]
One more word for describing what you lived and felt.
\end{grammarlens}

\subsection*{Your turn}
\begin{itemize}
\raggedright
  \item \emph{Say it:} \emph{upacarati}
  \item \emph{Say it:} \emph{upacarati}, once more
  \item \emph{Say it:} \emph{putrī rugṇāṃ mātaram upacarati.} — The daughter looks after her sick mother.
  \item \emph{Recall:} say the Sanskrit for a wholesome diet, then the Sanskrit for nursing, then say \emph{upacarati} again
\end{itemize}

\begin{tcolorbox}[breakable,colback=teal!4,colframe=teal!35!black,title={Before you move on}]
What does \sk{उपचरति} mean? (\textquotedblleft{}Attends on\textquotedblright{}.) Say it once more.
\end{tcolorbox}
\section[\texorpdfstring{\sk{आप्याययति} (āpyāyayati)}{आप्याययति (āpyāyayati)}]{\sk{आप्याययति} (āpyāyayati) — revives; refreshes}

\label{lesson:SA-C485-apyayayati}



\begin{quote}
Before the new one: say the Sanskrit for nursing, then the Sanskrit for attends on.
\end{quote}

\subsection*{You'll want to know: \sk{आप्याययति}}
\textbf{\sk{आप्याययति}} — \emph{āpyāyayati} — \textquotedblleft{}revives; refreshes\textquotedblright{}.

\begin{quote}
\sk{उष्णः} \sk{सूपः} \sk{आतुरम्} \sk{आप्याययति।}
\end{quote}

\begin{quote}
\emph{uṣṇaḥ sūpaḥ āturam āpyāyayati.}
\end{quote}

\begin{quote}
Warm soup refreshes the patient.
\end{quote}

Past (it …): \sk{सूपः} \sk{आतुरम्} \sk{आप्यायितवान्}, \emph{sūpaḥ āturam āpyāyitavān}.

\begin{cousinweb}
\textbf{\sk{आप्याययति}} (\emph{āpyāyayati}) is built from \emph{ā} and the root √pyai, \textquotedblleft{}to swell, to grow full\textquotedblright{}: to fill up again; MW's ppp \emph{āpyāyita} is \textquotedblleft{}satisfied, improved, pleased\textquotedblright{} (MW).
\end{cousinweb}

\begin{grammarlens}[title={What you've built}]
One more word for describing what you lived and felt.
\end{grammarlens}

\subsection*{Your turn}
\begin{itemize}
\raggedright
  \item \emph{Say it:} \emph{āpyāyayati}
  \item \emph{Say it:} \emph{āpyāyayati}, once more
  \item \emph{Say it:} \emph{uṣṇaḥ sūpaḥ āturam āpyāyayati.} — Warm soup refreshes the patient.
  \item \emph{Recall:} say the Sanskrit for nursing, then the Sanskrit for attends on, then say \emph{āpyāyayati} again
\end{itemize}

\begin{tcolorbox}[breakable,colback=teal!4,colframe=teal!35!black,title={Before you move on}]
What does \sk{आप्याययति} mean? (\textquotedblleft{}Revives; refreshes\textquotedblright{}.) Say it once more.
\end{tcolorbox}
\section[\texorpdfstring{\sk{उपशमः} (upaśamaḥ)}{उपशमः (upaśamaḥ)}]{\sk{उपशमः} (upaśamaḥ) — relief}

\label{lesson:SA-C485-upasamah}



\begin{quote}
Before the new one: say the Sanskrit for attends on, then the Sanskrit for revives.
\end{quote}

\subsection*{You'll want to know: \sk{उपशमः}}
\textbf{\sk{उपशमः}} — \emph{upaśamaḥ} — \textquotedblleft{}relief\textquotedblright{}.

\begin{quote}
\sk{औषधेन} \sk{वेदनायाः} \sk{उपशमः} \sk{भवति।}
\end{quote}

\begin{quote}
\emph{auṣadhena vedanāyāḥ upaśamaḥ bhavati.}
\end{quote}

\begin{quote}
The medicine brings relief from the pain.
\end{quote}

A masculine noun.

\begin{cousinweb}
\textbf{\sk{उपशमः}} (\emph{upaśamaḥ}) is built from \emph{upa} and the root √śam, \textquotedblleft{}to be calm\textquotedblright{}; MW gives \textquotedblleft{}assuagement, alleviation, cessation\textquotedblright{} (MW).
\end{cousinweb}

\begin{grammarlens}[title={What you've built}]
One more word for describing what you lived and felt.
\end{grammarlens}

\subsection*{Your turn}
\begin{itemize}
\raggedright
  \item \emph{Say it:} \emph{upaśamaḥ}
  \item \emph{Say it:} \emph{upaśamaḥ}, once more
  \item \emph{Say it:} \emph{auṣadhena vedanāyāḥ upaśamaḥ bhavati.} — The medicine brings relief from the pain.
  \item \emph{Recall:} say the Sanskrit for attends on, then the Sanskrit for revives, then say \emph{upaśamaḥ} again
\end{itemize}

\begin{tcolorbox}[breakable,colback=teal!4,colframe=teal!35!black,title={Before you move on}]
What does \sk{उपशमः} mean? (\textquotedblleft{}Relief\textquotedblright{}.) Say it once more.
\end{tcolorbox}
\section[\texorpdfstring{\sk{रसायनम्} (rasāyanam)}{रसायनम् (rasāyanam)}]{\sk{रसायनम्} (rasāyanam) — a rejuvenating tonic}

\label{lesson:SA-C485-rasayanam}



\begin{quote}
Before the new one: say the Sanskrit for revives, then the Sanskrit for relief.
\end{quote}

\subsection*{You'll want to know: \sk{रसायनम्}}
\textbf{\sk{रसायनम्}} — \emph{rasāyanam} — \textquotedblleft{}a rejuvenating tonic\textquotedblright{}.

\begin{quote}
\sk{वैद्यः} \sk{वृद्धाय} \sk{रसायनं} \sk{ददाति।}
\end{quote}

\begin{quote}
\emph{vaidyaḥ vṛddhāya rasāyanaṃ dadāti.}
\end{quote}

\begin{quote}
The physician gives the old man a tonic.
\end{quote}

A neuter noun.

\begin{cousinweb}
\textbf{\sk{रसायनम्}} (\emph{rasāyanam}) is built from \textbf{\sk{रसः}}, \textquotedblleft{}sap, essence\textquotedblright{}, and \emph{ayana}, \textquotedblleft{}a way\textquotedblright{}; MW gives \textquotedblleft{}a medicine supposed to prevent old age and prolong life, an elixir\textquotedblright{} (MW).
\end{cousinweb}

\begin{grammarlens}[title={What you've built}]
One more word for describing what you lived and felt.
\end{grammarlens}

\subsection*{Your turn}
\begin{itemize}
\raggedright
  \item \emph{Say it:} \emph{rasāyanam}
  \item \emph{Say it:} \emph{rasāyanam}, once more
  \item \emph{Say it:} \emph{vaidyaḥ vṛddhāya rasāyanaṃ dadāti.} — The physician gives the old man a tonic.
  \item \emph{Recall:} say the Sanskrit for revives, then the Sanskrit for relief, then say \emph{rasāyanam} again
\end{itemize}

\begin{tcolorbox}[breakable,colback=teal!4,colframe=teal!35!black,title={Before you move on}]
What does \sk{रसायनम्} mean? (\textquotedblleft{}A rejuvenating tonic\textquotedblright{}.) Say it once more.
\end{tcolorbox}
